\maketitle
\thispagestyle{plain}
\begin{abstract}

本文研究含光伏与储能的小区微网购电调度问题。四个问题使用同一套能量守恒与储能状态递推关系，真正变化的是\textbf{决策时刻能够获得的信息}以及\textbf{计划偏差的结算方式}。据此，本文先建立统一的母线侧调度模型，再依次处理确定性日前计划、历史信息下的风险购电、多时刻预报下的滚动调整以及波动电价下的重算，并始终按照“决策时只使用当时可获得的信息、执行后再更新状态”的原则推进。

\textbf{对于问题一}，将附件1给出的电价、负载和光伏预测视为确定信息，以全天购电费用最小为目标，对外网购电与储能充放电进行联合优化，并要求储能在日初、日末均保持 $6000$ kWh。为避免同一最优费用下出现无意义的充放电循环，采用“费用--储能吞吐量--弃光量”的词典序目标。结果表明，最优全天计划购电量为 \QoneEnergy\ kWh，购电费为 \QoneCost\ 元；与不配置储能时的 \QoneNoStorage\ 元相比，节省 \QoneSave\ 元，降幅为 \QoneSavePct\%。

\textbf{对于问题二}，题目没有给出当天的负载、光伏预报，因此主模型采用严格的历史信息口径：每日 $0{:}00$ 只利用此前已经发生的数据预测当天负载和光伏，再锁定全天计划购电量。考虑到预测光伏偏高会引致 $5$ 倍电价的紧急购电，单时段边际损失呈 $4{:}1$ 非对称，由此得到光伏低分位风险修正的解析启发；全年则在多个历史预测器和残差分位之间滚动选择。实际执行时，储能根据当期已实现偏差实时缓冲，仍无法覆盖的缺口才紧急购电。$2$ 月至 $12$ 月总费用为 \QtwoTotal\ 元，紧急购电量为 \QtwoEmergency\ kWh。作为信息边界对照，若提前知道当天完整负荷曲线，费用可降低 \DeltaLoad\ 元，说明负荷信息质量对全年成本具有显著影响。

\textbf{对于问题三}，在 $0{:}00$ 先形成全天基准计划，在 $6{:}00$、$12{:}00$、$18{:}00$ 获得新预报后，只对尚未执行的时段重新优化，并以当时实测储电量衔接前后两个滚动阶段。本文主模型按“计划购电费、调整相关费用和紧急购电费”叠加口径结算，并同时保留差额结算作为敏感性对照。仅在 $0{:}00$ 决策时，全年费用为 \QthreeSzero\ 元；依次加入三个日内更新节点后，费用降至 \QthreeSone、\QthreeStwo、\QthreeSthree\ 元，累计节省 \QthreeVtotal\ 元。进一步的配对对照表明，节点收益同时来自已实现负荷信息与新光伏预报，因此本文将其解释为\textbf{决策更新的综合价值}，而不把全部收益归因于光伏预报本身。

\textbf{对于问题四}，保持问题二、三的预测与滚动机制不变，将固定日电价替换为附件4给出的逐日实时电价重新求解。按题面直接使用附件4电价时，Q4-2 与 Q4-3 的全年费用分别为 \QfourTwoMain\ 元和 \QfourThreeMain\ 元，滚动更新节省 \QfourRollPerfect\ 元。另设置只利用历史价格的因果预测情形作为现实性扩展，结果表明价格信息受限会增加成本，但不改变“日内滚动更新能够降低费用”的主要结论。

全文对母线平衡、储能状态递推、容量与功率边界、费用回代和时间因果性进行独立校验，并对终端储电量与效率口径开展敏感性分析。结果说明，\textbf{储能能够在时间维度上转移低价电与光伏富余，而严格区分计划时可用信息与执行时实现信息，是避免调度结果过度乐观的关键。}

\keywords{微网经济调度；储能优化；因果预测；风险定价；滚动优化}
\end{abstract}
